\section*{Limitations}
\label{sec:limitations}
\paragraph{Scope of the prototype.}
The REST runtime is a dry run: it never writes kernel settings, and the separate actuator changes only an owned child thread's timer slack and affinity. The artifact therefore demonstrates a decision contract, a bounded per-thread actuator, raw measurements and non-LLM selector replay, but not a deployed VM-sysctl bundle actuator, traffic isolation or a verified trained checkpoint. We make no deployment-efficacy, production-safety or LLM-superiority claim; Section~\ref{sec:protocol} and Appendix~\ref{app:testbed} list the matched comparisons and deployment requirements such claims require. The staged adapter is not wired into the REST VM-sysctl runtime and does not establish traffic isolation or recovery from irreversible effects. The reasoner awaits model samples synchronously, so we make no fast-path latency or resource-footprint claim; a nonblocking deployment would separate synchronous enforcement from asynchronous proposals and version the objective weights, policy, graph and calibration snapshots per decision. Logs are append-only files within a process, without tamper resistance. Prompting a model to obey graph edges does not enforce those edges; independent validation is required, and the current REST runtime enforces only its explicit example constraints.

\paragraph{Task and annotation.}
The audited task has four candidate configurations, its evidence determines the best choice, and the default is best on the test grid; the audit can therefore expose missing added value and misreported evidence, but it cannot show that an LLM adds value on harder tasks. Claim labels were assigned by a single annotator (an author) without an agreement study; the per-claim labels and checks are released so that they can be re-annotated. The claim sample is small (97 claims), local explanations were byte-identical across seeds, and five seeds under greedy decoding are not independent samples.

\paragraph{Evaluation breadth.}
All controlled measurements use one periodic workload family on one shared host, with roles separated by period rather than by workload family or hardware. Measured graph induction on the training sweep emits no edge. Selector, BO and RL comparisons replay recorded runs rather than executing closed loops, and control is best on the tested grid, so these comparisons can reveal only losses. Run-level intervals describe variability on a shared host and do not account for uncontrolled interference, warm-cache bias or correlated system state; exploratory intervals are unadjusted for multiple comparisons and temporal dependence, and the closed-loop local load does not characterize queueing under an externally fixed arrival rate. The model audit is retrospective and small, and hosted runs are not reproducible bit for bit. The available telemetry prototype contains scheduler and block-I/O demonstration proxies, including timestamp-derived buckets; those are not end-to-end request latencies and must not calibrate an empirical SLO claim. The local benchmarks record completion, failure and deadline data for their narrow tasks but do not substitute for web serving, TPC-C, Kafka/Spark, GPU inference or audio pipelines.

\paragraph{Risk statements.}
The rank gate's assumptions can fail under drift, selective labeling, or a score model updated using calibration data, and the observed miss rates under a period shift already exceed the nominal level. Because calibration and test periods differ, runs share a host, and the proposition is marginal under exchangeability, the gate replay neither verifies its deployment assumptions nor establishes recovery. A canary can suffer harm before detection, and rollback cannot undo every side effect. No no-worse-than-default, catastrophic-regression prevention, Pareto superiority or universal safety claim follows from the prototype. Valid citations in model explanations establish provenance, not faithful or correct reasoning. These limitations define concrete measurements and engineering prerequisites for the next evaluation.
